\section{Banco de questões — Controle Externo}

\subsection*{N1 — Fundamento competitivo}
\begin{questao}{\nivel{N1} 06.01 — Cebraspe/Certo ou Errado}No plano federal, o controle externo é exercido pelo Congresso Nacional com auxílio do Tribunal de Contas da União.\end{questao}
\begin{questao}{\nivel{N1} 06.02 — Cebraspe/Certo ou Errado}O Tribunal de Contas integra o Poder Legislativo como órgão hierarquicamente subordinado, razão pela qual o Parlamento pode rever livremente suas decisões técnicas.\end{questao}
\begin{questao}{\nivel{N1} 06.03 — Cebraspe/Certo ou Errado}A fiscalização constitucional abrange legalidade, legitimidade e economicidade, não se limitando à regularidade formal da despesa.\end{questao}
\begin{questao}{\nivel{N1} 06.04 — Cebraspe/Certo ou Errado}No paradigma federal, o TCU julga as contas anuais do Presidente da República.\end{questao}
\begin{questao}{\nivel{N1} 06.05 — Cebraspe/Certo ou Errado}No paradigma federal, o TCU julga as contas dos administradores e demais responsáveis por dinheiros, bens e valores públicos sujeitos à sua jurisdição.\end{questao}
\begin{questao}{\nivel{N1} 06.06 — Cebraspe/Certo ou Errado}Uma irregularidade formal implica, necessariamente, existência de dano ao erário e obrigação de ressarcimento.\end{questao}
\begin{questao}{\nivel{N1} 06.07 — Cebraspe/Certo ou Errado}Controle interno e controle externo podem atuar de forma complementar, sem que um substitua integralmente o outro.\end{questao}
\begin{questao}{\nivel{N1} 06.08 — Cebraspe/Certo ou Errado}Denúncia apresentada ao Tribunal de Contas equivale, por si só, a prova da irregularidade narrada.\end{questao}
\begin{questao}{\nivel{N1} 06.09 — Cebraspe/Certo ou Errado}Medida cautelar possui finalidade preventiva e instrumental, não devendo ser confundida com sanção definitiva.\end{questao}
\begin{questao}{\nivel{N1} 06.10 — Cebraspe/Certo ou Errado}O monitoramento de deliberações deve verificar apenas se o gestor respondeu ao Tribunal, sendo desnecessário examinar implementação e efeito.\end{questao}

\subsection*{N2 — Aplicação}
\begin{questao}{\nivel{N2} 06.11 — FGV/contas}As contas anuais do chefe do Poder Executivo, no modelo constitucional, são objeto de:\begin{enumerate}[label=\Alph*)]\item parecer prévio do Tribunal de Contas e julgamento pelo Legislativo competente;\item julgamento definitivo exclusivo do Tribunal de Contas;\item sentença judicial;\item homologação do controle interno;\item arquivamento automático.\end{enumerate}\end{questao}
\begin{questao}{\nivel{N2} 06.12 — FCC/gestão}As contas de administrador que gerencia diretamente recursos públicos submetem-se, em regra, ao:\begin{enumerate}[label=\Alph*)]\item julgamento técnico do Tribunal de Contas competente;\item parecer prévio sem decisão;\item julgamento pelo chefe do Executivo;\item controle social exclusivamente;\item STF.\end{enumerate}\end{questao}
\begin{questao}{\nivel{N2} 06.13 — Cebraspe/Certo ou Errado}A apreciação das contas de governo e o julgamento de contas de gestão possuem objetos e efeitos institucionais distintos.\end{questao}
\begin{questao}{\nivel{N2} 06.14 — FGV/economicidade}Um órgão adquiriu equipamento pelo menor preço, mas com vida útil muito inferior e manutenção cara. A fiscalização de economicidade deve considerar:\begin{enumerate}[label=\Alph*)]\item apenas o preço de aquisição;\item custo e benefício global da solução, qualidade e condições de mercado;\item somente regularidade da nota fiscal;\item apenas autorização orçamentária;\item exclusivamente dolo do gestor.\end{enumerate}\end{questao}
\begin{questao}{\nivel{N2} 06.15 — FCC/legalidade-legitimidade}Despesa cumpre formalmente a lei, mas é realizada para finalidade estranha ao interesse público. O controle pode examinar, além da legalidade:\begin{enumerate}[label=\Alph*)]\item legitimidade;\item apenas contabilidade;\item somente competência penal;\item mérito partidário;\item nacionalidade do gestor.\end{enumerate}\end{questao}
\begin{questao}{\nivel{N2} 06.16 — Cebraspe/Certo ou Errado}Auditoria operacional pode examinar eficiência e efetividade, desde que haja critérios adequados e metodologia compatível com o objetivo.\end{questao}
\begin{questao}{\nivel{N2} 06.17 — FGV/débito}A equipe identifica falha procedimental sem dano financeiro demonstrado. A conclusão correta é:\begin{enumerate}[label=\Alph*)]\item débito é automático;\item irregularidade e dano são categorias distintas; é necessário demonstrar dano e responsabilidade para imputação de débito;\item toda falha gera solidariedade;\item somente o controle interno pode atuar;\item o processo deve ser arquivado.\end{enumerate}\end{questao}
\begin{questao}{\nivel{N2} 06.18 — FCC/responsabilidade}Para responsabilizar agente por dano, é relevante examinar:\begin{enumerate}[label=\Alph*)]\item conduta, nexo causal, dano e demais requisitos do regime aplicável;\item apenas cargo ocupado;\item somente valor do salário;\item popularidade do gestor;\item existência de denúncia anônima.\end{enumerate}\end{questao}
\begin{questao}{\nivel{N2} 06.19 — Cebraspe/Certo ou Errado}Responsabilidade solidária não deve ser presumida somente porque duas pessoas atuavam na mesma unidade administrativa.\end{questao}
\begin{questao}{\nivel{N2} 06.20 — FGV/cautelar}Licitação está prestes a ser homologada e há indício robusto de cláusula restritiva capaz de causar dano grave. Uma medida cautelar pode ser adequada para:\begin{enumerate}[label=\Alph*)]\item antecipar a multa definitiva;\item preservar a utilidade da decisão e evitar dano, se presentes os requisitos;\item condenar criminalmente os responsáveis;\item substituir toda a licitação por decisão do Tribunal;\item dispensar contraditório em qualquer fase definitivamente.\end{enumerate}\end{questao}
\begin{questao}{\nivel{N2} 06.21 — Cebraspe/Certo ou Errado}A adoção de cautelar não dispensa a posterior instrução e decisão de mérito segundo o devido processo.\end{questao}
\begin{questao}{\nivel{N2} 06.22 — FCC/controle interno}O controle interno deve, entre outras funções constitucionais, apoiar o controle externo e avaliar metas, legalidade e resultados nos termos do modelo aplicável.\begin{enumerate}[label=\Alph*)]\item correta;\item incorreta, porque só o Tribunal de Contas controla;\item incorreta, pois controle interno é privado;\item correta apenas em empresas;\item incorreta, porque não pode avaliar resultados.\end{enumerate}\end{questao}
\begin{questao}{\nivel{N2} 06.23 — FGV/representação}Uma representação aponta possível sobrepreço, acompanhada de documentos. O Tribunal deve:\begin{enumerate}[label=\Alph*)]\item considerar os fatos provados automaticamente;\item avaliar admissibilidade, instruir e formar convicção a partir das evidências;\item aplicar sanção sem defesa;\item transferir o julgamento ao denunciante;\item ignorar por não ser prestação de contas anual.\end{enumerate}\end{questao}
\begin{questao}{\nivel{N2} 06.24 — Cebraspe/Certo ou Errado}Controle concomitante pode permitir correção antes da consumação do dano, mas não deve transformar o órgão de controle em gestor da decisão administrativa.\end{questao}
\begin{questao}{\nivel{N2} 06.25 — FCC/monitoramento}Uma determinação exige revogar acessos de ex-servidores. No monitoramento, a melhor evidência é:\begin{enumerate}[label=\Alph*)]\item ofício dizendo ``cumprido'';\item relação atual de contas, registros de revogação e testes/amostras do processo preventivo;\item promessa verbal;\item manual de segurança;\item orçamento da ferramenta.\end{enumerate}\end{questao}

\subsection*{N3 — Integração de concurso}
\begin{questao}{\nivel{N3} 06.26 — FGV/assertivas}Considere: I. O controle externo pode avaliar economicidade. II. O Tribunal de Contas é órgão do Judiciário. III. O controle interno pode apoiar o externo. Assinale:\begin{enumerate}[label=\Alph*)]\item apenas I;\item apenas II;\item I e III;\item II e III;\item I, II e III.\end{enumerate}\end{questao}
\begin{questao}{\nivel{N3} 06.27 — Cebraspe/Certo ou Errado}O auxílio prestado pelo Tribunal de Contas ao Legislativo não elimina competências decisórias próprias atribuídas diretamente pela Constituição.\end{questao}
\begin{questao}{\nivel{N3} 06.28 — FCC/parecer prévio}O parecer prévio possui função de:\begin{enumerate}[label=\Alph*)]\item subsidiar o julgamento político das contas de governo pelo Legislativo competente;\item condenar criminalmente o chefe do Executivo;\item substituir a LOA;\item julgar definitivamente qualquer gestor;\item exercer controle interno.\end{enumerate}\end{questao}
\begin{questao}{\nivel{N3} 06.29 — FGV/ato sujeito a registro}A apreciação para fins de registro de determinados atos de pessoal é competência constitucional dos Tribunais de Contas, observadas:\begin{enumerate}[label=\Alph*)]\item apenas escolhas do gestor;\item regras constitucionais, legais e jurisprudenciais aplicáveis;\item votação popular;\item sanção penal;\item dispensa de processo.\end{enumerate}\end{questao}
\begin{questao}{\nivel{N3} 06.30 — Cebraspe/Certo ou Errado}O exame de legalidade de admissão e concessão de aposentadoria para registro não equivale a julgamento jurisdicional do Poder Judiciário.\end{questao}
\begin{questao}{\nivel{N3} 06.31 — FCC/tomada de contas especial}A tomada de contas especial é adequada, em linhas gerais, quando:\begin{enumerate}[label=\Alph*)]\item se pretende substituir qualquer prestação ordinária;\item há omissão/irregularidade/dano a apurar e os mecanismos ordinários não foram suficientes, conforme normas aplicáveis;\item não há responsável possível;\item se deseja aprovar previamente contrato;\item qualquer cidadão solicita informação.\end{enumerate}\end{questao}
\begin{questao}{\nivel{N3} 06.32 — FGV/processo}Responsável é sancionado sem oportunidade de defesa em processo de contas. O vício mais evidente relaciona-se a:\begin{enumerate}[label=\Alph*)]\item economicidade;\item contraditório e ampla defesa;\item orçamento;\item competência municipal;\item neutralidade.\end{enumerate}\end{questao}
\begin{questao}{\nivel{N3} 06.33 — Cebraspe/Certo ou Errado}O contraditório em processo de controle externo não impede a adoção de medidas cautelares em situações urgentes quando o ordenamento admitir contraditório diferido e estiverem presentes os requisitos.\end{questao}
\begin{questao}{\nivel{N3} 06.34 — FCC/sanção}A aplicação de multa pelo Tribunal de Contas depende de:\begin{enumerate}[label=\Alph*)]\item mera discordância técnica;\item competência e previsão legal, além de responsabilidade e processo adequado;\item aprovação do fornecedor;\item decisão judicial penal anterior;\item dano financeiro em qualquer hipótese.\end{enumerate}\end{questao}
\begin{questao}{\nivel{N3} 06.35 — FGV/determinação}Uma deliberação ordena que o órgão cumpra obrigação legal objetiva. O comando aproxima-se de:\begin{enumerate}[label=\Alph*)]\item recomendação meramente facultativa;\item determinação;\item parecer prévio;\item denúncia;\item consulta informal.\end{enumerate}\end{questao}
\begin{questao}{\nivel{N3} 06.36 — Cebraspe/Certo ou Errado}Recomendação pode buscar melhoria sem necessariamente impor a escolha de uma solução tecnológica específica quando várias alternativas são juridicamente adequadas.\end{questao}
\begin{questao}{\nivel{N3} 06.37 — FCC/prescrição}Em matéria sancionadora e de ressarcimento, o auditor deve:\begin{enumerate}[label=\Alph*)]\item presumir imprescritibilidade de tudo;\item identificar o regime constitucional, legal e jurisprudencial aplicável à pretensão examinada;\item ignorar o tempo decorrido;\item aplicar prazo civil automaticamente;\item considerar apenas a data do relatório.\end{enumerate}\end{questao}
\begin{questao}{\nivel{N3} 06.38 — FGV/controle operacional}Um programa cumpre todas as formalidades, mas entrega apenas 20\% da meta com custo crescente. A fiscalização mais adequada deve incluir:\begin{enumerate}[label=\Alph*)]\item apenas legalidade;\item desempenho, eficiência e efetividade, com critérios apropriados;\item somente processo penal;\item controle de constitucionalidade concentrado;\item parecer sobre eleição.\end{enumerate}\end{questao}
\begin{questao}{\nivel{N3} 06.39 — Cebraspe/Certo ou Errado}Uma auditoria de TI pode integrar o controle externo quando examina sistemas, dados, segurança ou contratações que afetam a gestão pública.\end{questao}
\begin{questao}{\nivel{N3} 06.40 — FCC/accountability}Prestação de contas, em sentido material, envolve:\begin{enumerate}[label=\Alph*)]\item somente enviar documentos;\item demonstrar e justificar gestão de recursos e resultados, sujeitando-se à avaliação e consequências;\item apenas publicar salário;\item cumprir auditoria privada;\item substituir responsabilidade do gestor.\end{enumerate}\end{questao}

\subsection*{N4 — Auditoria / avançado}
\begin{questao}{\nivel{N4} 06.41 — Cebraspe/caso integrado}Uma contratação é formalmente autorizada e empenhada, mas custa 70\% acima de soluções equivalentes sem justificativa. Julgue: a ausência de ilegalidade formal aparente não impede exame de economicidade pelo controle externo.\end{questao}
\begin{questao}{\nivel{N4} 06.42 — FGV/responsabilização}O relatório identifica dano de R\$ 1 milhão, mas não demonstra que o ato do gestor contribuiu causalmente para a perda. A imputação de débito:\begin{enumerate}[label=\Alph*)]\item é automática pela posição hierárquica;\item exige demonstração de responsabilidade e nexo conforme o regime aplicável;\item depende apenas do valor;\item deve ser solidária a todos os servidores;\item não exige defesa.\end{enumerate}\end{questao}
\begin{questao}{\nivel{N4} 06.43 — FCC/controle concomitante}Durante fiscalização, auditores começam a escolher especificações técnicas da licitação e autorizar cada alteração do edital. O principal risco institucional é:\begin{enumerate}[label=\Alph*)]\item controle insuficiente;\item assunção de função gerencial, comprometendo independência e responsabilidade;\item ausência de publicidade;\item excesso de contraditório;\item falta de sanção.\end{enumerate}\end{questao}
\begin{questao}{\nivel{N4} 06.44 — Cebraspe/Certo ou Errado}O controle preventivo deve reduzir risco sem transformar o Tribunal em coadministrador das escolhas fiscalizadas.\end{questao}
\begin{questao}{\nivel{N4} 06.45 — FGV/cautelar}Uma cautelar suspende provisoriamente ato até análise de mérito. Para diferenciá-la de sanção, é relevante notar que a cautelar:\begin{enumerate}[label=\Alph*)]\item possui finalidade instrumental/preventiva;\item declara culpa definitiva;\item exige dano consumado;\item encerra o processo;\item substitui a decisão de mérito.\end{enumerate}\end{questao}
\begin{questao}{\nivel{N4} 06.46 — Cebraspe/Certo ou Errado}O fato de o controle interno ter aprovado uma despesa não vincula o Tribunal de Contas nem impede fiscalização externa posterior.\end{questao}
\begin{questao}{\nivel{N4} 06.47 — FCC/denúncia}Uma denúncia apresenta alegação grave, mas sem qualquer suporte mínimo e em desacordo com requisitos de admissibilidade. O Tribunal deve:\begin{enumerate}[label=\Alph*)]\item condenar o gestor;\item aplicar as regras de admissibilidade e instrução, sem tratar alegação como prova;\item divulgar dados sigilosos;\item imputar débito automaticamente;\item dispensar decisão motivada.\end{enumerate}\end{questao}
\begin{questao}{\nivel{N4} 06.48 — FGV/monitoramento}O órgão cumpriu formalmente determinação para criar política de backup, mas testes mostram que restauração continua falhando. A conclusão é:\begin{enumerate}[label=\Alph*)]\item cumprimento integral e efetivo;\item implementação documental sem demonstração de efetividade do controle;\item inexistência de dever de monitorar;\item erro apenas de redação;\item necessidade de arquivar o processo.\end{enumerate}\end{questao}
\begin{questao}{\nivel{N4} 06.49 — Cebraspe/Certo ou Errado}Uma recomendação de auditoria deve preferencialmente tratar a causa/resultado esperado, preservando à gestão a escolha entre soluções lícitas quando não houver única alternativa necessária.\end{questao}
\begin{questao}{\nivel{N4} 06.50 — FGV/matriz de achado}Auditoria identifica pagamento por serviço sem evidência de medição, mas a contratada afirma verbalmente ter cumprido. O achado mais adequado é:\begin{enumerate}[label=\Alph*)]\item ``A contratada fraudou'';\item ``A ausência de evidência verificável de medição impede demonstrar a correspondência entre serviço entregue e pagamento, elevando risco de pagamento indevido; devem ser obtidos logs, registros de aceite e critérios contratuais'';\item ``Todo pagamento é irregular'';\item ``Basta solicitar nova nota fiscal'';\item ``A declaração verbal comprova o serviço''.\end{enumerate}\end{questao}
